\Askhsh{37139}{4}
\begin{ylh}{1}
    \item 3.4. 3ο Κριτήριο ισότητας τριγώνων
    \item 3.7. Κύκλος - Μεσοκάθετος - Διχοτόμος
    \item 3.11. Ανισοτικές σχέσεις πλευρών και γωνιών
    \item 4.2. Τέμνουσα δύο ευθειών - Ευκλείδειο αίτημα
    \item 5.6. Εφαρμογές στα τρίγωνα
    \item 5.11. Ισοσκελές τραπέζιο.
\end{ylh}
% ===================================================================
%  Αρχείο: 37139.tex (Fragment)
%  Αυτόματη μετατροπή μέσω doc2latex.py
% ===================================================================

ΘΕΜΑ 4

Δίνεται ισοσκελές τραπέζιο ΑΒΓΔ (ΑΒ//ΔΓ) και Ο το σημείο τομής των διαγωνίων του. Η ΑΓ είναι κάθετη στην ΑΔ και η ΒΔ είναι κάθετη στην ΒΓ. Θεωρούμε τα μέσα Μ, Ε και Ζ των ΓΔ, ΒΔ και ΑΓ αντίστοιχα.

Να αποδείξετε ότι:

\begin{alist}
\item ΜΕ=ΜΖ. \monades{6}

\item Η ΜΖ είναι κάθετη στην ΑΓ. \monades{6}

\item Τα τρίγωνα ΜΔΕ και ΜΖΓ είναι ίσα. \monades{7}

\item Η ΟΜ είναι μεσοκάθετος του ΕΖ. \monades{6}

\includesvg[width=5.5625in,height=2.73958in]{/home/spyros/Μαθηματικά/Trapeza_Thematwn/A_Lykeiou/Geometry/extracted/37139/media/image2.svg}
\end{alist}
